\subsubsection{Data Integration}
The Fortran module provides comprehensive routines for data integration, based on Jensen-Shannon Divergence. This is very useful for cross-study data.

\textbf{Key Features:}
\begin{itemize}
    \item \textbf{Calculate Jensen-Shannon Compatibility Score}: Score to evaluate how compatible two studies are.
    \item \textbf{Perform JSD Permutation Test}: Random sampling of the input data and recalculation of the JS score to test stability of the compatibility.
    \item \textbf{Family-Wise Score Contribution}: Calculates the impact of selected families on the calculated score.
\end{itemize}

\begin{enumerate}
    \item \textbf{Pool Mean Expression Values (Allocated)}

        Pools per-gene mean expression values from two studies, sorts them, and computes reference points \lstinline!x_star! using empirical quantiles.

        \textbf{Arguments:}
        \begin{itemize}
            \item \lstinline!n_genes_S1 integer(int32)!: number of genes in study S1
            \item \lstinline!mean_S1 real(real64), dimension(n_genes_S1)!: per-gene mean expression values for study S1
            \item \lstinline!n_genes_S2 integer(int32)!: number of genes in study S2
            \item \lstinline!mean_S2 real(real64), dimension(n_genes_S2)!: per-gene mean expression values for study S2
            \item \lstinline!n_points integer(int32)!: number of reference points to define
            \item \lstinline!n_pool integer(int32)!: number of included (non-NaN) pooled mean-expression values
            \item \lstinline!x_star real(real64), dimension(n_points)!: mean-expression reference points
            \item \lstinline!ierr integer(int32)!: error code
        \end{itemize}

        \begin{lstlisting}[caption={Pool Means (Allocated)}]
call pool_means_alloc(n_genes_S1, mean_S1, n_genes_S2, mean_S2, n_points, n_pool, x_star, ierr)
        \end{lstlisting}

    \item \textbf{Pool Mean Expression Values}

        Pools per-gene mean expression values across studies, removes trailing NaNs (which are sorted to the end), and computes empirical-quantile reference points \lstinline!x_star!.

        \textbf{Arguments:}
        \begin{itemize}
            \item \lstinline!pooled_means real(real64), dimension(pool_size)!: pooled mean-expression values
            \item \lstinline!pooled_means_perm integer(int32), dimension(pool_size)!: permutation vector that sorts \lstinline!pooled_means!
            \item \lstinline!pool_size integer(int32)!: number of pooled means (typically \lstinline!n_genes_S1 + n_genes_S2!)
            \item \lstinline!n_points integer(int32)!: number of reference points to compute
            \item \lstinline!n_pool integer(int32)!: number of included (non-NaN) pooled means
            \item \lstinline!x_star real(real64), dimension(n_points)!: mean-expression reference points
            \item \lstinline!ierr integer(int32)!: error code
        \end{itemize}

        \begin{lstlisting}[caption={Pool Means}]
call pool_means(pooled_means, pooled_means_perm, pool_size, n_points, n_pool, x_star, ierr)
        \end{lstlisting}

    \item \textbf{Calculate Neighborhood Size}

        Determines the number of neighbors to use for \lstinline!construct_neighborhoods!, based on the pooled mean-expression distribution, the number of reference points, and optional user-defined limits.

        \textbf{Arguments:}
        \begin{itemize}
            \item \lstinline!n_pool integer(int32)!: total number of pooled mean-expression values across both studies
            \item \lstinline!n_points integer(int32)!: number of reference points
            \item \lstinline!n_genes_S integer(int32)!: number of genes in the current study
            \item \lstinline!mean_S real(real64), dimension(n_genes_S)!: per-gene mean expression values
            \item \lstinline!desired_size integer(int32), optional!: optional upper limit for neighborhood size (default: 1000)
            \item \lstinline!n_neighbors integer(int32)!: calculated neighborhood size
        \end{itemize}

        \begin{lstlisting}[caption={Calculate Neighborhood Size}]
n_neighbors = calc_neighborhood_size(n_pool, n_points, n_genes_S, mean_S, desired_size)
        \end{lstlisting}

    \item \textbf{Construct Neighborhoods (Allocated)}

        Constructs neighborhood-based residual sets (kNN) by allocating temporary distance and permutation arrays, then delegating to \lstinline!construct_neighborhoods!.

        \textbf{Arguments:}
        \begin{itemize}
            \item \lstinline!n_points integer(int32)!: number of reference points
            \item \lstinline!x_star real(real64), dimension(n_points)!: mean-expression reference points
            \item \lstinline!n_genes_S integer(int32)!: number of genes in the study
            \item \lstinline!mean_S real(real64), dimension(n_genes_S)!: per-gene mean expression values
            \item \lstinline!n_reps_S integer(int32)!: number of biological replicates
            \item \lstinline!resid_S real(real64), dimension(n_reps_S, n_genes_S)!: matrix of signed residuals
            \item \lstinline!neighborhood_residuals real(real64), dimension(n_reps_S, n_neighbors, n_points)!: residual vectors for each neighborhood
            \item \lstinline!neighborhood_indices integer(int32), dimension(n_neighbors, n_points)!: indices of selected neighborhood genes
            \item \lstinline!n_neighbors integer(int32)!: number of neighbors (compute via \lstinline!calc_neighborhood_size!)
            \item \lstinline!ierr integer(int32)!: error code
        \end{itemize}

        \begin{lstlisting}[caption={Construct Neighborhoods (Allocated)}]
call construct_neighborhoods_alloc(n_points, x_star, n_genes_S, mean_S, n_reps_S, resid_S, &
    neighborhood_residuals, neighborhood_indices, n_neighbors, ierr)
        \end{lstlisting}

    \item \textbf{Construct Neighborhoods}

        Constructs neighborhood-based residual sets (kNN) for each reference point by computing distances to \lstinline!x_star!, sorting genes by distance, and selecting the \lstinline!n_neighbors! closest non‑NaN genes.

        \textbf{Arguments:}
        \begin{itemize}
            \item \lstinline!n_points integer(int32)!: number of reference points
            \item \lstinline!x_star real(real64), dimension(n_points)!: mean‑expression reference points
            \item \lstinline!n_genes_S integer(int32)!: number of genes in the study
            \item \lstinline!mean_S real(real64), dimension(n_genes_S)!: per‑gene mean expression values
            \item \lstinline!n_reps_S integer(int32)!: number of biological replicates
            \item \lstinline!resid_S real(real64), dimension(n_reps_S, n_genes_S)!: matrix of signed residuals
            \item \lstinline!tmp_distances real(real64), dimension(n_genes_S)!: work array for distances
            \item \lstinline!tmp_distances_perm integer(int32), dimension(n_genes_S)!: permutation vector for sorting distances
            \item \lstinline!neighborhood_residuals real(real64), dimension(n_reps_S, n_neighbors, n_points)!: residual vectors for each neighborhood
            \item \lstinline!neighborhood_indices integer(int32), dimension(n_neighbors, n_points)!: indices of selected neighborhood genes
            \item \lstinline!n_neighbors integer(int32)!: number of neighbors (compute via \lstinline!calc_neighborhood_size!)
            \item \lstinline!ierr integer(int32)!: error code
        \end{itemize}

        \begin{lstlisting}[caption={Construct Neighborhoods}]
call construct_neighborhoods(n_points, x_star, n_genes_S, mean_S, n_reps_S, resid_S, &
    tmp_distances, tmp_distances_perm, neighborhood_residuals, neighborhood_indices, &
    n_neighbors, ierr)
        \end{lstlisting}

    \item \textbf{Determine Shared Residual Range}

        Computes the shared residual range \([-R, R]\) for the absolute residuals pooled from studies S1 and S2.

        \textbf{Arguments:}
        \begin{itemize}
            \item \lstinline!abs_residual_pool real(real64), dimension(pool_size)!: absolute residual values of the concatenated S1,S2 residuals
            \item \lstinline!abs_residual_pool_perm integer(int32), dimension(pool_size)!: permutation vector that sorts \lstinline!abs_residual_pool!
            \item \lstinline!pool_size integer(int32)!: size of the residual pool
            \item \lstinline!shared_residual_range real(real64)!: computed residual range \(R\)
            \item \lstinline!ierr integer(int32)!: error code
            \item \lstinline!residual_range_quantile real(real64), optional!: quantile for determining the residual range (default: 95.0)
        \end{itemize}

\begin{lstlisting}[caption={Shared Residual Range}]
call determine_shared_residual_range(abs_residual_pool, abs_residual_pool_perm, pool_size, shared_residual_range, ierr, residual_range_quantile)
\end{lstlisting}
    
    \item \textbf{Determine Shared Residual Range (Allocated)}

        Computes the shared residual range \([-R, R]\) for two studies by allocating and flattening the residual pool from both studies before computing the percentile-based range.

        \textbf{Arguments:}
        \begin{itemize}
            \item \lstinline!neighborhood_residuals_S1 real(real64), dimension(n_reps_S1, n_neighbors, n_points)!: neighborhood residuals for study 1, NaN allowed
            \item \lstinline!neighborhood_residuals_S2 real(real64), dimension(n_reps_S2, n_neighbors, n_points)!: neighborhood residuals for study 2, NaN allowed
            \item \lstinline!n_reps_S1 integer(int32)!: number of replicates in study 1
            \item \lstinline!n_reps_S2 integer(int32)!: number of replicates in study 2
            \item \lstinline!n_neighbors integer(int32)!: number of neighbors
            \item \lstinline!n_points integer(int32)!: number of reference points
            \item \lstinline!shared_residual_range real(real64)!: computed residual range \(R\)
            \item \lstinline!ierr integer(int32)!: error code
            \item \lstinline!residual_range_quantile real(real64), optional!: quantile for determining the residual range (default: 95.0)
        \end{itemize}

        \begin{lstlisting}[caption={Allocated Shared Residual Range}]
call determine_shared_residual_range_alloc(neighborhood_residuals_S1, neighborhood_residuals_S2, &
    n_reps_S1, n_reps_S2, n_neighbors, n_points, shared_residual_range, ierr, residual_range_quantile)
        \end{lstlisting}

    \item \textbf{Build Residual Histograms}

        Summarizes neighborhood residuals into absolute histogram counts and probability mass functions \lstinline!pmf(residual, bin)! for each reference point.

        \textbf{Arguments:}
        \begin{itemize}
            \item \lstinline!neighborhood_residuals real(real64), dimension(n_reps, n_neighbors, n_points)!: neighborhood residuals for a study, NaN allowed
            \item \lstinline!n_reps integer(int32)!: number of replicates
            \item \lstinline!n_neighbors integer(int32)!: number of neighbors
            \item \lstinline!n_points integer(int32)!: number of reference points
            \item \lstinline!shared_residual_range real(real64)!: residual range \(R\) from \lstinline!determine_shared_residual_range_alloc!
            \item \lstinline!n_bins integer(int32)!: number of histogram bins in \([-R, R]\)
            \item \lstinline!counts integer(int32), dimension(n_points, n_bins)!: absolute histogram counts
            \item \lstinline!pmf real(real64), dimension(n_points, n_bins)!: normalized histogram counts
            \item \lstinline!included_n_reps integer(int32), dimension(n_points)!: number of non-NaN residuals per reference point
            \item \lstinline!ierr integer(int32)!: error code
            \item \lstinline!neighbor_mask logical, dimension(n_neighbors, n_points), optional!: mask to exclude specific neighbors
        \end{itemize}

        \begin{lstlisting}[caption={Residual Histograms}]
call build_residual_histograms(neighborhood_residuals, n_reps, n_neighbors, n_points, &
    shared_residual_range, n_bins, counts, pmf, included_n_reps, ierr, neighbor_mask)
        \end{lstlisting}

    \item \textbf{Compute Divergence per Reference Point}

        Computes the Jensen–Shannon divergence per reference point using the probability mass functions from two studies.

        \textbf{Arguments:}
        \begin{itemize}
            \item \lstinline!pmf_S1 real(real64), dimension(n_points, n_bins)!: normalized histogram counts for study 1
            \item \lstinline!pmf_S2 real(real64), dimension(n_points, n_bins)!: normalized histogram counts for study 2
            \item \lstinline!n_points integer(int32)!: number of reference points
            \item \lstinline!n_bins integer(int32)!: number of histogram bins in \([-R, R]\)
            \item \lstinline!js_divergences real(real64), dimension(n_points)!: Jensen–Shannon divergence per reference point
            \item \lstinline!ierr integer(int32)!: error code
        \end{itemize}

        \begin{lstlisting}[caption={Per-Reference-Point Divergence}]
call compute_divergence_per_reference_point(pmf_S1, pmf_S2, n_points, n_bins, js_divergences, ierr)
        \end{lstlisting}

    \item \textbf{Compute Weighted Global Divergence}

        Computes the global weighted Jensen–Shannon divergence using per-reference-point divergences and the number of included residuals from both studies.

        \textbf{Arguments:}
        \begin{itemize}
            \item \lstinline!js_divergences real(real64), dimension(n_points)!: Jensen–Shannon divergence per reference point
            \item \lstinline!n_points integer(int32)!: number of reference points
            \item \lstinline!included_n_reps_S1 integer(int32), dimension(n_points)!: non-NaN residual counts for study 1
            \item \lstinline!included_n_reps_S2 integer(int32), dimension(n_points)!: non-NaN residual counts for study 2
            \item \lstinline!global_js_divergence real(real64)!: weighted global Jensen–Shannon divergence
            \item \lstinline!weights real(real64), dimension(n_points)!: weights used in the global divergence calculation
            \item \lstinline!ierr integer(int32)!: error code
        \end{itemize}

        \begin{lstlisting}[caption={Weighted Global Divergence}]
call compute_weighted_global_divergence(js_divergences, n_points, included_n_reps_S1, &
    included_n_reps_S2, global_js_divergence, weights, ierr)
        \end{lstlisting}

    \item \textbf{GJCT Permutation Test (Allocated)}

        Estimates how likely the observed global Jensen–Shannon divergence is to occur by chance under the null hypothesis that both studies are exchangeable.  
        Allocates all required working arrays and delegates to \lstinline!gjct_permutation_test!.

        \textbf{Arguments:}
        \begin{itemize}
            \item \lstinline!neighborhood_residuals_S1 real(real64), dimension(n_reps_S1, n_neighbors, n_points)!: neighborhood residuals for study 1, NaN allowed
            \item \lstinline!neighborhood_residuals_S2 real(real64), dimension(n_reps_S2, n_neighbors, n_points)!: neighborhood residuals for study 2, NaN allowed
            \item \lstinline!n_reps_S1 integer(int32)!: number of replicates in study 1
            \item \lstinline!n_reps_S2 integer(int32)!: number of replicates in study 2
            \item \lstinline!n_neighbors integer(int32)!: number of neighbors
            \item \lstinline!n_points integer(int32)!: number of reference points
            \item \lstinline!global_jsd_observed real(real64)!: observed global JSD value
            \item \lstinline!n_bins integer(int32)!: number of histogram bins
            \item \lstinline!shared_residual_range real(real64)!: shared residual range from \lstinline!determine_shared_residual_range!
            \item \lstinline!n_permutations integer(int32)!: number of permutations to perform
            \item \lstinline!jsd_null real(real64), dimension(n_permutations)!: global JSD values under the null hypothesis
            \item \lstinline!p_value real(real64)!: empirical p‑value
            \item \lstinline!ierr integer(int32)!: error code
            \item \lstinline!random_seed integer(int32), optional!: seed for shuffling
            \item \lstinline!neighbor_mask_S1 logical, dimension(n_neighbors, n_points), optional!: mask to exclude neighbors in study 1
            \item \lstinline!neighbor_mask_S2 logical, dimension(n_neighbors, n_points), optional!: mask to exclude neighbors in study 2
        \end{itemize}

        \begin{lstlisting}[caption={GJCT Permutation Test (Allocated)}]
call gjct_permutation_test_alloc(neighborhood_residuals_S1, neighborhood_residuals_S2, &
    n_reps_S1, n_reps_S2, n_neighbors, n_points, global_jsd_observed, n_bins, &
    shared_residual_range, n_permutations, jsd_null, p_value, ierr, random_seed, &
    neighbor_mask_S1, neighbor_mask_S2)
        \end{lstlisting}

    \item \textbf{GJCT Permutation Test}

        Estimates how likely the observed global Jensen–Shannon divergence is to occur by chance under the null hypothesis that both studies are exchangeable.  
        Reuses preallocated workspace arrays and performs all permutations in-place.

        \textbf{Arguments:}
        \begin{itemize}
            \item \lstinline!neighborhood_residuals_S1_copy real(real64), dimension(n_reps_S1, n_neighbors, n_points), intent(inout)!: copy of residuals for study 1, shuffled in-place
            \item \lstinline!neighborhood_residuals_S2_copy real(real64), dimension(n_reps_S2, n_neighbors, n_points), intent(inout)!: copy of residuals for study 2, shuffled in-place
            \item \lstinline!n_reps_S1 integer(int32)!: number of replicates in study 1
            \item \lstinline!n_reps_S2 integer(int32)!: number of replicates in study 2
            \item \lstinline!n_neighbors integer(int32)!: number of neighbors
            \item \lstinline!n_points integer(int32)!: number of reference points
            \item \lstinline!global_jsd_observed real(real64)!: observed global JSD value
            \item \lstinline!n_bins integer(int32)!: number of histogram bins
            \item \lstinline!shared_residual_range real(real64)!: shared residual range from \lstinline!determine_shared_residual_range!
            \item \lstinline!n_permutations integer(int32)!: number of permutations to perform
            \item \lstinline!jsd_null real(real64), dimension(n_permutations)!: global JSD values under the null hypothesis
            \item \lstinline!p_value real(real64)!: empirical p‑value
            \item \lstinline!tmp_pool real(real64), dimension(n_reps_S1 + n_reps_S2, n_neighbors)!: working array for shuffling pooled residuals
            \item \lstinline!tmp_pmf_S1 real(real64), dimension(n_points, n_bins)!: PMF workspace for study 1
            \item \lstinline!tmp_pmf_S2 real(real64), dimension(n_points, n_bins)!: PMF workspace for study 2
            \item \lstinline!tmp_counts integer(int32), dimension(n_points, n_bins)!: histogram count workspace
            \item \lstinline!tmp_included_n_reps_S1 integer(int32), dimension(n_points)!: workspace for included counts (study 1)
            \item \lstinline!tmp_included_n_reps_S2 integer(int32), dimension(n_points)!: workspace for included counts (study 2)
            \item \lstinline!tmp_js_divergences real(real64), dimension(n_points)!: workspace for per‑point divergences
            \item \lstinline!tmp_weights real(real64), dimension(n_points)!: workspace for weights
            \item \lstinline!ierr integer(int32)!: error code
            \item \lstinline!random_seed integer(int32), optional!: seed for shuffling
            \item \lstinline!neighbor_mask_S1 logical, dimension(n_neighbors, n_points), optional!: mask to exclude neighbors in study 1
            \item \lstinline!neighbor_mask_S2 logical, dimension(n_neighbors, n_points), optional!: mask to exclude neighbors in study 2
        \end{itemize}

        \begin{lstlisting}[caption={GJCT Permutation Test}]
call gjct_permutation_test(neighborhood_residuals_S1_copy, neighborhood_residuals_S2_copy, &
    n_reps_S1, n_reps_S2, n_neighbors, n_points, global_jsd_observed, n_bins, &
    shared_residual_range, n_permutations, jsd_null, p_value, tmp_pool, tmp_pmf_S1, &
    tmp_pmf_S2, tmp_counts, tmp_included_n_reps_S1, tmp_included_n_reps_S2, &
    tmp_js_divergences, tmp_weights, ierr, random_seed, neighbor_mask_S1, neighbor_mask_S2)
        \end{lstlisting}

    \item \textbf{FJCT Compute JSD (Allocated)}

        Computes the family-level compatibility score \lstinline!global_js_divergence! between two studies for a single gene family by reusing the same conditioning-on-mean-expression pipeline as the global gJCT, but restricting residual samples to genes belonging to the specified family.

        \textbf{Arguments:}
        \begin{itemize}
            \item \lstinline!family_idx integer(int32)!: index of the family to analyze
            \item \lstinline!gene_to_family_S1 integer(int32), dimension(n_genes_S1)!: mapping of genes to families in study 1
            \item \lstinline!gene_to_family_S2 integer(int32), dimension(n_genes_S2)!: mapping of genes to families in study 2
            \item \lstinline!n_genes_S1 integer(int32)!: number of genes in study 1
            \item \lstinline!n_genes_S2 integer(int32)!: number of genes in study 2
            \item \lstinline!n_reps_S1 integer(int32)!: number of replicates in study 1
            \item \lstinline!n_reps_S2 integer(int32)!: number of replicates in study 2
            \item \lstinline!n_neighbors integer(int32)!: number of neighbors
            \item \lstinline!n_points integer(int32)!: number of reference points
            \item \lstinline!neighborhood_residuals_S1 real(real64), dimension(n_reps_S1, n_neighbors, n_points)!: residuals for study 1, NaN allowed
            \item \lstinline!neighborhood_residuals_S2 real(real64), dimension(n_reps_S2, n_neighbors, n_points)!: residuals for study 2, NaN allowed
            \item \lstinline!neighborhood_genes_S1 integer(int32), dimension(n_reps_S1, n_neighbors, n_points)!: gene indices for study 1
            \item \lstinline!neighborhood_genes_S2 integer(int32), dimension(n_reps_S2, n_neighbors, n_points)!: gene indices for study 2
            \item \lstinline!n_bins integer(int32)!: number of histogram bins
            \item \lstinline!shared_residual_range real(real64)!: shared residual range from \lstinline!determine_shared_residual_range!
            \item \lstinline!js_divergences real(real64), dimension(n_points)!: Jensen–Shannon divergence per reference point
            \item \lstinline!included_n_reps_S1 integer(int32), dimension(n_points)!: included non-NaN residuals for study 1
            \item \lstinline!included_n_reps_S2 integer(int32), dimension(n_points)!: included non-NaN residuals for study 2
            \item \lstinline!total_included_n_reps integer(int32)!: total included residuals across both studies
            \item \lstinline!global_js_divergence real(real64)!: weighted global Jensen–Shannon divergence
            \item \lstinline!weights real(real64), dimension(n_points)!: weights used in the global divergence calculation
            \item \lstinline!ierr integer(int32)!: error code
        \end{itemize}

        \begin{lstlisting}[caption={Family-Level JSD (Allocated)}]
call fjct_compute_jsd_alloc(family_idx, gene_to_family_S1, gene_to_family_S2, &
    n_genes_S1, n_genes_S2, neighborhood_residuals_S1, neighborhood_residuals_S2, &
    neighborhood_genes_S1, neighborhood_genes_S2, n_reps_S1, n_reps_S2, n_neighbors, &
    n_points, n_bins, shared_residual_range, js_divergences, included_n_reps_S1, &
    included_n_reps_S2, total_included_n_reps, global_js_divergence, weights, ierr)
        \end{lstlisting}

    \item \textbf{FJCT Compute JSD}

        Computes the family-level compatibility score \lstinline!global_js_divergence! between two studies for a single gene family, reusing preallocated workspace arrays for efficiency.

        \textbf{Arguments:}
        \begin{itemize}
            \item \lstinline!family_idx integer(int32)!: index of the family to analyze
            \item \lstinline!gene_to_family_S1 integer(int32), dimension(n_genes_S1)!: mapping of genes to families in study 1
            \item \lstinline!gene_to_family_S2 integer(int32), dimension(n_genes_S2)!: mapping of genes to families in study 2
            \item \lstinline!n_genes_S1 integer(int32)!: number of genes in study 1
            \item \lstinline!n_genes_S2 integer(int32)!: number of genes in study 2
            \item \lstinline!n_reps_S1 integer(int32)!: number of replicates in study 1
            \item \lstinline!n_reps_S2 integer(int32)!: number of replicates in study 2
            \item \lstinline!n_neighbors integer(int32)!: number of neighbors
            \item \lstinline!n_points integer(int32)!: number of reference points
            \item \lstinline!neighborhood_residuals_S1 real(real64), dimension(n_reps_S1, n_neighbors, n_points)!: residuals for study 1, NaN allowed
            \item \lstinline!neighborhood_residuals_S2 real(real64), dimension(n_reps_S2, n_neighbors, n_points)!: residuals for study 2, NaN allowed
            \item \lstinline!neighborhood_genes_S1 integer(int32), dimension(n_reps_S1, n_neighbors, n_points)!: gene indices for study 1
            \item \lstinline!neighborhood_genes_S2 integer(int32), dimension(n_reps_S2, n_neighbors, n_points)!: gene indices for study 2
            \item \lstinline!n_bins integer(int32)!: number of histogram bins
            \item \lstinline!shared_residual_range real(real64)!: shared residual range from \lstinline!determine_shared_residual_range!
            \item \lstinline!js_divergences real(real64), dimension(n_points)!: Jensen–Shannon divergence per reference point
            \item \lstinline!included_n_reps_S1 integer(int32), dimension(n_points)!: included non-NaN residuals for study 1
            \item \lstinline!included_n_reps_S2 integer(int32), dimension(n_points)!: included non-NaN residuals for study 2
            \item \lstinline!total_included_n_reps integer(int32)!: total included residuals across both studies
            \item \lstinline!global_js_divergence real(real64)!: weighted global Jensen–Shannon divergence
            \item \lstinline!weights real(real64), dimension(n_points)!: weights used in the global divergence calculation
            \item \lstinline!tmp_counts integer(int32), dimension(n_points, n_bins)!: working array for histogram construction
            \item \lstinline!pmf_S1 real(real64), dimension(n_points, n_bins)!: PMF for study 1
            \item \lstinline!pmf_S2 real(real64), dimension(n_points, n_bins)!: PMF for study 2
            \item \lstinline!ierr integer(int32)!: error code
        \end{itemize}

        \begin{lstlisting}[caption={Family-Level JSD}]
call fjct_compute_jsd(family_idx, gene_to_family_S1, gene_to_family_S2, &
    n_genes_S1, n_genes_S2, neighborhood_residuals_S1, neighborhood_residuals_S2, &
    neighborhood_genes_S1, neighborhood_genes_S2, n_reps_S1, n_reps_S2, n_neighbors, &
    n_points, n_bins, shared_residual_range, js_divergences, included_n_reps_S1, &
    included_n_reps_S2, total_included_n_reps, global_js_divergence, weights, &
    tmp_counts, pmf_S1, pmf_S2, ierr)
        \end{lstlisting}
\end{enumerate}
